% Compile locally with Tectonic or XeLaTeX. Segoe UI headings, Calibri body.
\documentclass[letterpaper,10pt]{article}
\usepackage{resume_style}
\hypersetup{pdftitle={Yuhao Xie | LLM Inference and GPU Systems}}
\begin{document}
\fontsize{10.3}{12.5}\selectfont

\cvhead{LLM Inference \enspace\textbar\enspace KV Cache Quantization \enspace\textbar\enspace GPU Kernel Optimization}{%
Computer Engineering PhD student focused on efficient LLM inference, with an \textbf{EMNLP 2026 main-conference first-author paper} and experience integrating CUDA quantization kernels into an LLM inference engine.}

\section{EDUCATION}
\compactentry{University of Central Florida}{PhD, Computer Engineering}{Aug. 2023 -- Aug. 2027 (expected)}
\sep
\compactentry{University of Florida}{MS, Electrical \& Computer Engineering}{Aug. 2021 -- May 2023}
\sep
\compactentry{Hangzhou Dianzi University}{BS, Electrical \& Computer Engineering}{Sept. 2017 -- June 2021}

\section{PUBLICATION}
\cvpublication{\textbf{Yuhao Xie}, Mingjie Lin.}
{Beyond: Better-than-Full-Precision KV Caches via Learnable Non-Uniform Quantization as Implicit Regularization.}
{EMNLP 2026, main conference.}
{\href{https://wsxhjnb1.github.io/Beyond-KV-Cache/paper/Beyond.pdf}{\textcolor{accent}{[Paper]}}
\enspace\href{https://github.com/wsxhjnb1/Beyond-KV-Cache}{\textcolor{accent}{[Code]}}
\enspace\href{https://wsxhjnb1.github.io/Beyond-KV-Cache/}{\textcolor{accent}{[Project]}}}

\section{LLM RESEARCH \& ENGINEERING}
\entry{Hithink RoyalFlush Information Network}{May -- Aug. 2024}{Research Intern, Large Language Models}{Zhejiang, China}
\begin{itemize}
  \item Developed CUDA quantization/dequantization kernels, reducing execution time by \textbf{90\% versus a PyTorch implementation}; integrated them into the company's LLM inference engine with \textbf{FlashAttention and paged KV-cache support}.
  \item Implemented \textbf{2-bit KV-cache quantization} with separate grouping along the sequence and head dimensions for keys and values; achieved \textbf{50\%} lower GPU memory use and \textbf{8$\times$} throughput for the company's 7B financial LLMs \mbox{on a single NVIDIA A10}.
  \item Tuned quantization group size and residual length on internal benchmarks to retain \textbf{99.3\% of the original model's accuracy}.
\end{itemize}

\section{SELECTED SYSTEMS PROJECTS}
\project{Fused Attention Decoding with Quantized KV Caches}
\begin{itemize}
  \item Fused \textbf{on-the-fly KV dequantization} with attention scoring, online softmax, and value accumulation; reconstructed packed K/V in registers using groupwise quantization parameters, avoiding full-precision KV materialization in HBM.
  \item Achieved a \textbf{40\% inference speedup for Llama 2 7B} on a single \textbf{NVIDIA H100} versus a Hugging Face implementation.
\end{itemize}
\sep
\project{OrbitCall: vLLM Agent Scoring Optimization \enspace\href{https://github.com/wsxhjnb1/OrbitCall}{\textcolor{accent}{[Code]}}}
\begin{itemize}
  \item Built a \textbf{vLLM 0.29 suffix-logprob plugin for CUDA GPUs} to score finite key-order tool-call families; preserved native \textbf{LM-head projection shapes} and complete JSON/EOS likelihoods while reducing normalization, rank calculation, and CPU transfer.
  \item On \textbf{DGX Spark (GB10)}, measured a \textbf{1.08$\times$ complete two-round Agent speedup} for Qwen2.5-3B versus native full scoring across six paired blocks (eight fixed workflows).
\end{itemize}

\section{ADDITIONAL ENGINEERING EXPERIENCE}
\compactentry{Rtelligent Technology}{Software Engineer Intern}{Jan. -- May 2021}
\begin{itemize}
  \item Developed ARM Cortex-M firmware in C++; reduced G-code processing time by \textbf{20\%} and used Xilinx FPGA raster-image parallelism to increase laser engraving speed by \textbf{30\%}.
\end{itemize}
\section{TECHNICAL SKILLS}
\textbf{Programming \& tools:} Python, C/C++, CUDA, OpenMP, Verilog, Git, Docker.\par
\textbf{ML systems:} PyTorch, Transformers, vLLM, FlashAttention, KV-cache quantization, speculative decoding.\par

\begingroup
\titlespacing*{\section}{0pt}{5pt}{4pt}
\section{TEACHING, SERVICE \& HONORS}
\fontsize{10.3}{11.8}\selectfont
\textbf{UCF: Instructor:} Aug. -- Dec. 2026: Digital Systems (EEE3342C) and \mbox{Computer Organization (EEL3801C)}; \mbox{Aug. 2024 -- May 2025: Computer Organization Lab}; \textbf{Teaching Assistant:} Aug. -- Dec. 2026: Hardware Security and Trusted Circuit Design (EEE4346C); \mbox{May -- Aug. 2025: Junior Design}.\par
\textbf{Reviewer:} ICLR 2027.\par
\textbf{ORCGS Doctoral Fellowship}, UCF (2023--2027).\par
Zhejiang provincial collegiate programming competition, team third prize (2019); academic scholarships, Hangzhou Dianzi University (2017--2020).
\par\endgroup

\end{document}
